\begin{adjustwidth}{-1cm}{-1cm}
%% Abstract (configurado para língua inglesa)
\begin{resumo}[Abstract]
\begin{otherlanguage*}{english}
Academic performance in Brazilian public high schools is strongly conditioned by sociological factors such as poverty, youth labor, digital exclusion, and age-grade distortion, yet access to individual student records is restricted by legal and institutional barriers. This work proposes a predictive model of academic performance trends for students of the state public school network of Piauí, grounded in sociological factors and trained on a synthetic population representative of the profile of an average class of that network. Every generation parameter is traceable to official statistics and academic literature, and the generator incorporates anti-deterministic mechanisms: student financial assistance is derived from a legal eligibility rule, vulnerabilities interact non-linearly, and academic resilience profiles are preserved. Observable academic indicators correspond to first-semester partial records, and the target variable represents the end-of-year trend, labeled according to Brazilian educational legislation on approval criteria. Over three thousand records, three classification algorithms and two reference \textit{baselines} were compared under stratified cross-validation, with oversampling confined to each training partition and paired statistical testing. XGBoost achieved 84.11\% accuracy, statistically equivalent to \textit{Random Forest} and superior to the single tree and to the \textit{baselines}, with no errors between the extreme classes. After the partial academic signals, the most determinant attributes were age-grade distortion, financial assistance, youth labor, and per capita income, recovering the relative weight of the sociological factors in the generating phenomenon. It is concluded that the articulation between traceable synthetic data and ensemble classifiers constitutes a valid methodological alternative for early-warning systems in restricted-data contexts.

\vspace{\onelineskip}
\noindent
\textbf{Keywords}: academic performance; sociological factors; synthetic data; machine learning; public high school.
\end{otherlanguage*}
\end{resumo}
\end{adjustwidth}
